\paragraph{Task Interdependence in Task-based Models of Automation.}
Our framework relates to the literature on task-based models of technical change, which generally view tasks as predetermined and technologically independent objects combined through a CES production function \citep{acemoglu2011skills, acemoglu2018, acemoglu2019automation, acemoglu2024tasks}.
We depart from this literature by modeling production as a sequence of technologically interdependent steps, where AI chaining creates complementarities across adjacent steps that can overturn comparative advantage as the driver of who does what.\footnote{Two contemporary works by \cite{trammell2026workflows} and \cite{farboodi2026data} study a different form of linkage across tasks, where production still occurs task by task, but performing a task raises productivity on others through learning, by workers in the former and by machines in the latter. Our emphasis is instead on the ordering of tasks in production.}
Chaining also lets step-level decisions aggregate into the system-level substitution that \cite{autor2003skill}, \cite{acemoglu2019automation}, and \cite{bresnahan2002information} emphasize, since entire chains can be automated in a single leap rather than step by step.


\paragraph{Complementarities, O-Ring Dynamics, and the Productivity J-Curve.}
Our work contributes to the literature on friction-laden technology adoption and organizational complementarities, in which information technologies raise productivity mainly alongside complementary organizational change \citep{bresnahan2002information, bloom2016management}.
We micro-found these complementarities through the O-ring theory of production \citep{kremer1993oring}. Chaining makes the payoff to automating a step depend on whether its neighbors can join the same AI block, so the returns to improving any one step turn on performance elsewhere and arrive only once AI quality crosses a threshold that makes longer chains worthwhile.
Our results align with \citet{gans2026oring}, who likewise emphasizes task interdependence in O-ring production and obtains lumpy automation through time reallocation rather than workflow adjacency.
Because these returns arrive only once AI quality passes such a threshold, the marginal benefit of improving AI is non-monotone, and in this sense our framework provides a micro-foundation for the ``productivity J-curve'' pattern of technology adoption \citep{brynjolfsson2021productivity, mcelheran2025rise, furman2019ai, tambe2012productivity, bonney2024tracking, mcelheran2024adoption}.


\paragraph{Division of Labor and the Boundary of the Job.}
Our work also connects to the literature on the division of labor within a firm.
The Coase--Williamson logic, in which economic boundaries are shaped by coordination, transaction, and governance costs \citep{coase1937nature, williamson1971vertical, williamson1979transaction}, applies to the ``boundary of the job'' as well, in that the firm chooses how many and which tasks to bundle into a worker's role, making the degree of specialization endogenous \citep{murphy1986specialization, becker1992division, dessein2006adaptive}.
Our hand-off costs are the intra-firm analogue of the transaction costs that AI cannot fully subsume, capturing the tacit or social knowledge aspect of work \citep{deming2017growing}, and playing the role that the cost of unbundling tasks plays in \citet{garicano2026weak}. 
AI nonetheless redraws job boundaries by stripping skill out of the steps it absorbs, effectively changing the skill requirements of jobs \citep{autor2024new, autor2025expertise, autor2026makes}.\footnote{
\cite{ide2025ai} provides a distinct but related view, embedding AI as an autonomous problem solver in a knowledge hierarchy \citep{garicano2000hierarchies, garicano2006organization} and studying how it reshapes spans of control within the firm.
}


\paragraph{Measuring AI Exposure and Realized Execution.}
A rapidly growing literature has sought to quantify the impact of new automation technologies on the labor market.
Early work studies the susceptibility of occupations to computerization \citep{frey2017future}, and more recent work develops measures of occupational exposure to AI \citep{brynjolfsson2018can, webb2020impact, felten2021occupational, eloundou2023gpts, hampole2025artificial, tomlinson2025working}, typically by mapping technologies to specific tasks and then aggregating these task-level mappings to construct occupation-level indices.
We build on these measures by showing how, in addition to the share of tasks exposed to AI, the arrangement of those tasks in a workflow also governs the gains from automation.